% =============================================================================
%  TEMPLATE USP REFINADO v1.0
%  Dissertação/Tese — Universidade de São Paulo
%  Baseado em abntex2, conforme diretrizes SIBi-USP
%  Compilar: pdflatex -> pdflatex (2x)
% =============================================================================
\documentclass[
    % -- opções para dissertação/tese --
    12pt,                   % tamanho da fonte
    a4paper,                % tamanho do papel
    openright,              % capítulos começam em página ímpar
    oneside,                % impressão em frente e verso
    english,                % idioma adicional para resumo
    brazil                  % idioma principal
]{abntex2}

% --- Pacotes fundamentais ---
\usepackage[utf8]{inputenc}
\usepackage[T1]{fontenc}
\usepackage{amsmath,amssymb,amsthm}
\usepackage{graphicx}
\usepackage{xcolor}
\usepackage{microtype}
\usepackage{booktabs}
\usepackage{caption}
\usepackage{subcaption}
\usepackage{enumitem}
\usepackage{algorithm}
\usepackage{algpseudocode}
\usepackage{times}          % Times New Roman (padrão USP)

% --- Hiperlinks (ABNT exige preto para versão impressa) ---
% abntex2 já carrega hyperref internamente
\hypersetup{
    colorlinks=true,
    linkcolor=black,
    citecolor=black,
    urlcolor=black,
    breaklinks=true,
    bookmarksnumbered=true
}

% --- Citações ABNT (alfabética) ---
\usepackage[alf]{abntex2cite}

% --- Configuração de margens (ABNT NBR 14724/USP SIBi) ---
\setlrmarginsandblock{3cm}{2cm}{*}
\setulmarginsandblock{3cm}{2cm}{*}

% --- Ambiente teoremas ---
\theoremstyle{plain}
\newtheorem{teorema}{Teorema}[chapter]
\newtheorem{lema}[teorema]{Lema}
\newtheorem{corolario}[teorema]{Corol\'ario}
\theoremstyle{definition}
\newtheorem{definicao}[teorema]{Defini\c{c}\~ao}
\theoremstyle{remark}
\newtheorem{observacao}[teorema]{Observa\c{c}\~ao}

% --- Comandos auxiliares ---
\newcommand{\R}{\mathbb{R}}
\newcommand{\N}{\mathbb{N}}
\newcommand{\E}{\mathbb{E}}

% =============================================================================
\begin{document}

% --- Elementos pré-textuais ---

% Capa (obrigatório USP)
\renewcommand{\imprimircapa}{%
  \begin{capa}%
    \center
    % \includegraphics[width=3cm]{usp_logo.pdf}\\[2cm]
    \textbf{UNIVERSIDADE DE S\~AO PAULO}\\
    \textbf{Nome do Programa de P\'os-Gradua\c{c}\~ao}\\[1cm]
    \textbf{NOME COMPLETO DO AUTOR}\\[4cm]
    \textbf{\large T\'ITULO DA DISSERTA\c{C}\~AO OU TESE:\\
            SUBT\'ITULO QUANDO HOUVER}\\[3cm]
    \vfill
    S\~ao Paulo\\
    2026
  \end{capa}
}

% Folha de rosto (obrigatório USP)
\renewcommand{\imprimirfolhaderosto}{%
  \begin{folhaderosto}
    \center
    \textbf{NOME COMPLETO DO AUTOR}\\[3cm]
    \textbf{\large T\'ITULO DA DISSERTA\c{C}\~AO OU TESE:\\
            SUBT\'ITULO QUANDO HOUVER}\\[2cm]
    \hspace*{\fill}\begin{minipage}{8cm}
    \small
    Disserta\c{c}\~ao apresentada ao Programa de
    P\'os-Gradua\c{c}\~ao em [Nome do Programa] da
    Universidade de S\~ao Paulo para obten\c{c}\~ao do
    t\'itulo de Mestre em [\'Area].
    \end{minipage}\\[2cm]
    \vfill
    Orientador: Prof. Dr. Nome do Orientador\\[1cm]
    S\~ao Paulo\\
    2026
  \end{folhaderosto}
}

% --- Capa ---
\imprimircapa

% --- Folha de rosto ---
\imprimirfolhaderosto

% --- Ficha catalográfica (opcional, SIBi-USP gera automaticamente) ---
\begin{fichacatalografica}
    \vspace*{\fill}
    \begin{center}
    \fbox{\begin{minipage}{13cm}
        \small
        \textbf{FICHA CATALOGR\'AFICA}\\[0.5cm]
        Elaborada pelo Sistema de Bibliotecas da USP\\
        com os dados fornecidos pelo autor.
    \end{minipage}}
    \end{center}
\end{fichacatalografica}

% --- Folha de aprovação ---
\begin{folhadeaprovacao}
    \center
    \textbf{NOME COMPLETO DO AUTOR}\\[2cm]
    \textbf{T\'ITULO DA DISSERTA\c{C}\~AO OU TESE}\\[2cm]
    \hspace*{\fill}\begin{minipage}{8cm}
    \small Disserta\c{c}\~ao apresentada ao Programa de
    P\'os-Gradua\c{c}\~ao em [Nome] para obten\c{c}\~ao do
    t\'itulo de Mestre.
    \end{minipage}\\[2cm]
    \vfill
    \textbf{Banca Examinadora}\\[1cm]
    \rule{8cm}{0.5pt} \\ Prof. Dr. Presidente (USP) \\
    \rule{8cm}{0.5pt} \\ Prof. Dr. Titular 1 \\
    \rule{8cm}{0.5pt} \\ Prof. Dr. Titular 2 \\
    \vfill
    S\~ao Paulo, \_\_\_ de \_\_\_\_\_\_\_\_\_ de 2026.
\end{folhadeaprovacao}

% --- Dedicatória ---
\begin{dedicatoria}
    \vspace*{\fill}
    \centering
    \noindent
    \textit{Aos meus pais, pelo apoio incondicional.}
    \vspace*{\fill}
\end{dedicatoria}

% --- Agradecimentos ---
\begin{agradecimentos}
    Agrade\c{c}o ao meu orientador, Prof. Dr. Nome, pela
    orienta\c{c}\~ao e confian\c{c}a ao longo deste trabalho.
    \`A CAPES e ao CNPq, pelo suporte financeiro.
    Aos colegas de laborat\'orio, pelas discuss\~oes produtivas.
\end{agradecimentos}

% --- Epígrafe ---
\begin{epigrafe}
    \vspace*{\fill}
    \begin{flushright}
        \textit{``A ci\^encia \'e o conhecimento de como\\
        as coisas funcionam.''\\
        Richard Feynman}
    \end{flushright}
    \vspace*{\fill}
\end{epigrafe}

% --- Resumo em português ---
\begin{resumo}
    Este trabalho apresenta um template refinado para
    disserta\c{c}\~oes e teses da Universidade de S\~ao Paulo
    (USP), elaborado conforme as diretrizes do Sistema
    Integrado de Bibliotecas (SIBi-USP) e da ABNT.
    O template inclui todos os elementos pr\'e-textuais,
    textuais e p\'os-textuais exigidos pelo regulamento.

    \textbf{Palavras-chave}: USP; abntex2; template LaTeX;
    disserta\c{c}\~ao; tese.
\end{resumo}

% --- Abstract em inglês ---
\begin{abstract}
    This work presents a refined template for dissertations
    and theses at the University of S\~ao Paulo (USP),
    prepared according to the guidelines of the Integrated
    Library System (SIBi-USP) and ABNT standards.
    The template includes all pre-textual, textual, and
    post-textual elements required by the regulations.

    \textbf{Keywords}: USP; abntex2; LaTeX template;
    dissertation; thesis.
\end{abstract}

% --- Lista de figuras ---
\listoffigures*

% --- Lista de tabelas ---
\listoftables*

% --- Lista de abreviaturas e siglas ---
\begin{siglas}
    \item[ABNT] Associa\c{c}\~ao Brasileira de Normas T\'ecnicas
    \item[CAPES] Coordena\c{c}\~ao de Aperfei\c{c}oamento de
        Pessoal de N\'ivel Superior
    \item[CNPq] Conselho Nacional de Desenvolvimento
        Cient\'ifico e Tecnol\'ogico
    \item[SIBi] Sistema Integrado de Bibliotecas
    \item[USP] Universidade de S\~ao Paulo
\end{siglas}

% --- Sumário ---
\tableofcontents*

% =============================================================================
% --- Elementos textuais ---

\textual

% =============================================================================
\chapter{Introdu\c{c}\~ao}
\label{chap:intro}

A Universidade de S\~ao Paulo (USP) exige que todas as
disserta\c{c}\~oes e teses sigam normas espec\'ificas de
formata\c{c}\~ao, definidas pelo Sistema Integrado de
Bibliotecas (SIBi-USP) em conformidade com a ABNT.

Este template oferece uma implementa\c{c}\~ao completa e
refinada dessas exig\^encias, permitindo que o pesquisador
se concentre no conte\'udo cient\'ifico.

\section{Objetivos}

O objetivo geral deste trabalho \'e fornecer um ponto de
partida robusto para a elabora\c{c}\~ao de disserta\c{c}\~oes
e teses USP em \LaTeX.

\subsection{Objetivos Espec\'ificos}

\begin{enumerate}
    \item Implementar todos os elementos pr\'e-textuais
        exigidos pelo SIBi-USP;
    \item Oferecer estrutura completa para cap\'itulos,
        se\c{c}\~oes e subse\c{c}\~oes;
    \item Incluir ambientes formais (teoremas, lemas,
        algoritmos);
    \item Fornecer exemplos de cita\c{c}\~oes no estilo
        ABNT (alfab\'etico).
\end{enumerate}

% =============================================================================
\chapter{Fundamenta\c{c}\~ao Te\'orica}
\label{chap:fund}

\section{Formula\c{c}\~ao Formal}

Seja $\mathcal{X} \subseteq \R^n$ um espa\c{c}o de
configura\c{c}\~oes e $f: \mathcal{X} \to \R$ uma fun\c{c}\~ao
objetivo. O problema geral de otimiza\c{c}\~ao pode ser
definido como:

\begin{equation}
    x^* = \arg\min_{x \in \mathcal{X}} f(x)
    \quad \text{s.a.} \quad g_i(x) \leq 0,\;
    i = 1, \dots, m
    \label{eq:otim}
\end{equation}

\begin{teorema}[Converg\^encia do Gradiente Descendente]
\label{thm:conv}
Sob a hip\'otese de $f$ convexa e $L$-suave, o gradiente
descendente com passo $\eta < 2/L$ satisfaz:
\begin{equation}
    f(x_k) - f(x^*) \leq \frac{\|x_0 - x^*\|^2}{2\eta k}
    \label{eq:conv}
\end{equation}
\end{teorema}

\begin{proof}
A demonstra\c{c}\~ao segue da $L$-suavidade e da atualiza\c{c}\~ao
$ x_{k+1} = x_k - \eta \nabla f(x_k)$ \cite{Boyd2004Convex}.
\end{proof}

\section{Resultados Conhecidos}

A Tabela~\ref{tab:comparativo} apresenta resultados
comparativos da literatura.

\begin{table}[htbp]
\centering
\caption{Comparativo de m\'etodos de otimiza\c{c}\~ao}
\label{tab:comparativo}
\begin{tabular}{lccc}
\toprule
M\'etodo & Taxa & Mem\'oria & Garantia \\
\midrule
GD & $O(1/k)$ & $O(n)$ & Convexo \\
AGD & $O(1/k^2)$ & $O(n)$ & Convexo \\
Newton & Quadr\'atica & $O(n^2)$ & Fort. convexo \\
\bottomrule
\end{tabular}
\end{table}

% =============================================================================
\chapter{Metodologia}
\label{chap:metod}

A metodologia proposta \'e composta por tr\^es etapas
principais, descritas no Algoritmo~\ref{alg:metod}.

\begin{algorithm}[htbp]
\caption{Metodologia proposta}
\label{alg:metod}
\begin{algorithmic}[1]
\Require Dataset $\mathcal{D}$, par\^ametros $\Theta$
\Ensure Resultados $\mathcal{R}$
\State $\mathcal{D}' \gets \text{Pr\'e-processamento}(\mathcal{D})$
\State $\mathcal{M} \gets \text{Treinamento}(\mathcal{D}', \Theta)$
\State $\mathcal{R} \gets \text{Avalia\c{c}\~ao}(\mathcal{M}, \mathcal{D}')$
\State \Return $\mathcal{R}$
\end{algorithmic}
\end{algorithm}

% =============================================================================
\chapter{Resultados e Discuss\~ao}
\label{chap:result}

A Figura~\ref{fig:resultados} ilustra os resultados obtidos
com a metodologia proposta.

\begin{figure}[htbp]
\centering
\vspace{2cm}
\begin{minipage}{0.7\textwidth}
\centering
\fbox{\parbox{\textwidth}{\centering
    [Gr\'afico de resultados experimentais]}}
\end{minipage}
\caption{Resultados experimentais comparativos}
\label{fig:resultados}
\end{figure}

% =============================================================================
\chapter{Conclus\~ao}
\label{chap:conc}

Este template fornece uma base completa e refinada para
disserta\c{c}\~oes e teses USP, abrangendo todos os
elementos exigidos pelo SIBi-USP e ABNT.

\section*{Trabalhos Futuros}

Sugere-se como trabalhos futuros:
\begin{enumerate}
    \item Adapta\c{c}\~ao para diferentes programas de
        p\'os-gradua\c{c}\~ao da USP;
    \item Integra\c{c}\~ao com ferramentas de CI/CD para
        compila\c{c}\~ao automatizada;
    \item Gera\c{c}\~ao autom\'atica de ficha catalogr\'afica.
\end{enumerate}

% =============================================================================
% --- Elementos pós-textuais ---

\postextual

% --- Referências (ABNT alfabético) ---
\bibliography{usp_bibliografia}

% --- Apêndices ---
\begin{apendicesenv}
    \chapter{C\'odigo-Fonte Complementar}
    \label{ap:codigo}

    Este ap\^endice cont\'em o c\'odigo-fonte dos
    experimentos realizados.

    \chapter{Resultados Adicionais}
    \label{ap:res_adicionais}

    Resultados complementares n\~ao inclu\'idos no corpo
    principal do texto.
\end{apendicesenv}

% --- Anexos ---
\begin{anexosenv}
    \chapter{Parecer do Comit\^e de \'Etica}
    \label{an:etica}

    Documento de aprova\c{c}\~ao do Comit\^e de
    \'Etica em Pesquisa.
\end{anexosenv}

% --- Índice remissivo ---
\printindex

\end{document}
